\documentclass[11pt]{article}

\usepackage[margin=1in]{geometry}
\usepackage{graphicx}
\usepackage{amsmath}
\usepackage{listings}
\usepackage{xcolor}
\usepackage{hyperref}
\usepackage{enumitem}
\usepackage{booktabs}

\title{Improving FAISS IVF Search using CLIO CTE\\
Current Implementation, Observations and Future Work}

\author{Carmen Vázquez}
\date{\today}

\begin{document}

\maketitle

\tableofcontents

\newpage

%%%%%%%%%%%%%%%%%%%%%%%%%%%%%%%%%%%%%%%%%%%%%%%%%%%%%%%%%%%%%%%
\section{Introduction}

During the last few weeks I have been working on integrating
FAISS IVF search into the CLIO runtime using the Context Transfer
Engine (CTE). The original motivation was to improve the behaviour
of very large FAISS indexes that no longer fit in RAM.

This document explains the work that has been implemented so far,
the problem that motivated the project, the current implementation,
the results and limitations that I have observed during the
experiments and, finally, the change that I believe would allow CLIO
to remove the remaining bottleneck.

%%%%%%%%%%%%%%%%%%%%%%%%%%%%%%%%%%%%%%%%%%%%%%%%%%%%%%%%%%%%%%%
\section{Background}

\subsection{How FAISS IVF Works}

FAISS is a library for Approximate Nearest Neighbour (ANN) search.
Instead of comparing a query against every vector in the database,
an IVF (Inverted File Index) first partitions the dataset into
multiple clusters, usually called \textbf{inverted lists}.

Each vector belongs to exactly one inverted list.

During a search:

\begin{enumerate}
\item the query is assigned to its nearest clusters,
\item only those clusters are visited,
\item the vectors inside those clusters are compared against the query.
\end{enumerate}

This reduces the amount of work compared to scanning the complete
dataset.

Two parameters are particularly important:

\begin{itemize}
\item \textbf{nlist}: the number of inverted lists in the index.
\item \textbf{nprobe}: the number of lists visited during each search.
\end{itemize}

Increasing \texttt{nprobe} generally improves recall, because more
candidate vectors are examined, but it also increases the amount of
data that must be read during the search.

%%%%%%%%%%%%%%%%%%%%%%%%%%%%%%%%%%%%%%%%%%%%%%%%%%%%%%%%%%%%%%%
\subsection{FAISS OnDiskInvertedLists}

When the inverted lists become too large to fit into memory,
FAISS provides the OnDiskInvertedLists implementation.

Instead of storing every inverted list in RAM, all the lists are
stored inside a single \texttt{.ivfdata} file. This file is memory
mapped using the Linux \texttt{mmap()} system call, which allows the
operating system to transparently load data into memory only when it
is actually accessed.

This approach works well while the index still fits comfortably
inside RAM. The operating system keeps recently used pages in the
page cache, which is the region of memory where the kernel keeps
copies of file data it has already read, and future accesses are
usually served directly from memory.

However, this behaviour changes once the index becomes several
times larger than the available RAM.

%%%%%%%%%%%%%%%%%%%%%%%%%%%%%%%%%%%%%%%%%%%%%%%%%%%%%%%%%%%%%%%
\subsection{The Problem with mmap}

The operating system manages the mapped file through the page cache.

The important point is that the page cache does not understand what
an inverted list is. Instead, it only manages independent memory
pages of 4 kilobytes each.

Whenever a query needs an inverted list, the operating system loads
the required pages from disk. If memory is already full, other pages
must first be evicted, that is, removed from memory to make room.

When the index is much larger than RAM, this process repeats
continuously:

\begin{itemize}
\item pages are loaded,
\item older pages are evicted,
\item future queries need them again,
\item the same pages are loaded once more.
\end{itemize}

This behaviour is commonly known as \textbf{page cache thrashing}.

One of the clearest symptoms observed during the experiments is a
large increase in the number of \textbf{page-ins}. A page-in is a
single event where the operating system fetches one page of data
from disk into memory. Instead of reusing data already resident in
memory, the operating system repeatedly fetches the same pages from
disk.

As memory pressure increases, throughput drops significantly, and
the main cost of the search becomes waiting for data to be loaded
back into memory rather than computing distances.

%%%%%%%%%%%%%%%%%%%%%%%%%%%%%%%%%%%%%%%%%%%%%%%%%%%%%%%%%%%%%%%
\subsection{Motivation}

The previous observations suggest that the operating system is
making memory placement decisions without understanding the access
pattern of FAISS.

From the application's point of view, the natural unit of data is an
inverted list. From the operating system's point of view, however,
the only unit is an independent memory page.

The main motivation of this project was therefore to replace the
implicit page-cache management performed by the operating system with
explicit placement managed by CLIO's Context Transfer Engine.

Instead of relying on the kernel to decide which pages remain in
memory, the application should explicitly manage complete inverted
lists. This idea led to the implementation described in the following
sections.

\newpage

%%%%%%%%%%%%%%%%%%%%%%%%%%%%%%%%%%%%%%%%%%%%%%%%%%%%%%%%%%%%%%%
\section{Using CLIO CTE instead of mmap}

The main idea behind this project is to replace the operating
system's implicit page management with explicit placement managed by
CLIO. Instead of storing all inverted lists inside a single
memory-mapped file, every inverted list is stored independently
inside the Context Transfer Engine (CTE), as a separate object called
a \textbf{blob}.

This allows CLIO to treat each inverted list as an independent object
that can be placed, migrated or evicted on its own. Because the
inverted list is already the natural unit of access for FAISS,
storing one inverted list per CTE blob means that the unit CLIO
manages is exactly the unit the search reads.

CTE stores these blobs across two storage tiers: a capped RAM tier,
which is fast but limited in size, and a larger NVMe tier on the
local solid state disk, which is slower but can hold data that does
not fit in the RAM tier. CLIO decides, per blob, where each inverted
list lives.

%%%%%%%%%%%%%%%%%%%%%%%%%%%%%%%%%%%%%%%%%%%%%%%%%%%%%%%%%%%%%%%
\subsection{Architecture}

The implementation consists of three main components:

\begin{itemize}
\item A preprocessing tool that converts a FAISS IVF index into CTE blobs.
\item A ChiMod, which is a module running inside the CLIO runtime, that
performs the IVF search.
\item The Context Transfer Engine, which stores the inverted lists
across the available storage tiers.
\end{itemize}

The overall architecture is shown conceptually below.

\begin{verbatim}

             Search request

        +----------------------+
        |       Client         |
        +----------+-----------+
                   |
                   |
          Shared Memory Task
                   |
                   v
        +----------------------+
        |    FAISS ChiMod      |
        |                      |
        | Quantization         |
        | List Selection       |
        | Distance Computation |
        +----------+-----------+
                   |
         Read required lists
                   |
                   v
        +----------------------+
        |         CTE          |
        |                      |
        | Blob 0               |
        | Blob 1               |
        | Blob 2               |
        | ...                  |
        +----------------------+

\end{verbatim}

Unlike the original FAISS implementation, the inverted lists are
never stored inside a single large file. Instead, every list becomes
an independent CTE blob.

%%%%%%%%%%%%%%%%%%%%%%%%%%%%%%%%%%%%%%%%%%%%%%%%%%%%%%%%%%%%%%%
\subsection{Index Ingestion}

Before any search can be performed, an existing FAISS IVF index must
be imported into CTE.

A standalone ingestion tool reads the populated IVF index, extracts
every inverted list and writes it into its corresponding CTE blob.
The ingestion process performs the following steps:

\begin{enumerate}
\item Load the FAISS IVF index.
\item Read the size of every inverted list.
\item Create one CTE blob for each non-empty list.
\item Copy the list contents into that blob.
\item Finally, store a \texttt{sizes} blob.
\end{enumerate}

Each list blob stores the encoded vectors followed by their
identifiers, packed together. The \texttt{sizes} blob is a small
array with one entry per list, giving the number of vectors in each
list. It acts as a directory: from it the search knows which lists
exist, how many bytes each blob contains, and where inside a blob the
vectors end and the identifiers begin. It is written last, so its
presence also signals that the ingestion finished correctly.

The result is a CTE representation where every inverted list can be
managed independently.

%%%%%%%%%%%%%%%%%%%%%%%%%%%%%%%%%%%%%%%%%%%%%%%%%%%%%%%%%%%%%%%
\subsection{Opening an Index}

The ChiMod only loads the metadata required to perform the search. In
particular, it loads:

\begin{itemize}
\item the coarse quantizer, which is the small structure used to assign
each query to its nearest clusters,
\item the IVF metadata (the number of lists, the vector dimension and
the code size),
\item the \texttt{sizes} blob,
\item the total number of vectors.
\end{itemize}

The inverted lists themselves are \textbf{not} loaded during this
step. Instead, they remain inside CTE and are fetched only when a
search actually needs them.

This keeps the memory footprint small when opening very large
indexes, and it is what allows an index larger than RAM to be opened
at all.

%%%%%%%%%%%%%%%%%%%%%%%%%%%%%%%%%%%%%%%%%%%%%%%%%%%%%%%%%%%%%%%
\subsection{Search Execution}

Every search follows the same sequence of operations.

First, the coarse quantizer assigns every query to its nearest
clusters. Those assignments determine which inverted lists must be
visited.

Since many queries often access the same lists, the implementation
builds the union of all required lists, so that each list is read at
most once even when several queries need it.

For every required list, the ChiMod issues an asynchronous
\texttt{AsyncGetBlob()} request to CTE. Several requests are kept in
flight at the same time in order to hide part of the storage latency,
so that while some lists are still being read, others are already
being processed. Whenever one read completes, the corresponding
inverted list is scanned immediately: the vectors inside it are
compared against every query that requires it, updating each query's
Top-K results, that is, the list of the K nearest neighbours found so
far for that query. To accelerate this step, the scan of each list is
parallelized across several worker threads.

Finally, after every required list has been processed, the Top-K
results are returned to the client.

%%%%%%%%%%%%%%%%%%%%%%%%%%%%%%%%%%%%%%%%%%%%%%%%%%%%%%%%%%%%%%%
\subsection{Expected Benefits}

Compared with the original mmap-based implementation, this design
moves the responsibility for data placement from the operating system
page cache to CTE. Because each inverted list is stored independently,
placement and eviction happen at exactly the same granularity that the
search algorithm uses, and different storage tiers can be managed
explicitly without modifying FAISS itself.

At this point, the implementation successfully replaces
\texttt{OnDiskInvertedLists} with CTE-based storage and performs
correct IVF search inside the CLIO runtime. The following section
discusses the experimental results and the remaining performance
bottleneck discovered during the evaluation.

%%%%%%%%%%%%%%%%%%%%%%%%%%%%%%%%%%%%%%%%%%%%%%%%%%%%%%%%%%%%%%%
\section{Experimental Observations}

After implementing the CTE-based storage, I first validated that the
ChiMod produced exactly the same results as the original FAISS
implementation, and then measured its behaviour on indexes of
increasing size.

\subsection{Setup}

All the experiments were run on a single compute node with 48
gigabytes of RAM and a local NVMe solid state disk. Each measurement
uses 500 queries, K equal to 10, \texttt{nprobe} equal to
\texttt{nlist} divided by 64, and 8 threads. Every configuration is
run for three passes: one \textbf{cold} pass, immediately after
evicting the data from memory, and two \textbf{warm} passes.

The main metric is \textbf{QPS} (queries per second), that is, how
many search queries the system can answer each second. A higher value
is better. As a side metric I also record the number of page-ins per
pass, which indicates how much data the operating system had to fetch
from disk.

\subsection{Correctness}

Before measuring performance, I checked correctness. For every
configuration, the distances and identifiers returned by the ChiMod
are \textbf{bitwise-identical} to those returned by stock FAISS. In
other words, the new storage backend does not change the result of
the search at all, only the way the data is stored and accessed.

\subsection{Results}

Table~\ref{tab:qps} shows the throughput of the CTE-based
implementation as the index size grows well beyond the size of the
RAM tier.

\begin{table}[h]
\centering
\begin{tabular}{lcccc}
\toprule
Index size & nprobe & QPS (cold) & QPS (warm) & page-ins per pass \\
\midrule
4.8 GB  & 64  & 148 & 145 - 152 & 0 \\
24 GB   & 128 & 21  & 27 - 30   & 0 \\
49 GB   & 256 & 8.9 & 7.8 - 7.9 & 0 \\
87 GB   & 256 & 4.0 & 3.9 - 4.0 & 0 \\
\bottomrule
\end{tabular}
\caption{CTE-based IVF search: throughput versus index size.}
\label{tab:qps}
\end{table}

Two facts stand out immediately. First, the number of page-ins is
zero at every size: the search never triggers the operating system to
fetch pages from disk, because CTE, not the kernel, manages the data.
Second, the cold and warm passes give almost the same throughput. The
system does not need a warm-up phase, and its performance is
predictable.

This is very different from the original memory-mapped baseline
measured on the same node. With \texttt{OnDiskInvertedLists}, the cold
pass of the 24 gigabyte index runs at about 5.7 QPS with more than two
hundred thousand page-ins, and only reaches roughly 70 QPS once the
data is warm in the page cache. For indexes several times larger than
RAM, the memory-mapped baseline degrades much further, dropping to a
small fraction of one query per second with continuous page-cache
thrashing. The CTE-based implementation, in contrast, keeps a stable
throughput of several queries per second at 49 and 87 gigabytes, with
no page faults and no collapse.

This confirms that the original goal was achieved: replacing the page
cache with explicit CTE placement removes the page-cache thrashing and
the performance collapse in the regime where the index is much larger
than RAM.

\subsection{A New Bottleneck}

Although this solves the original problem, the same experiments reveal
a different limitation.

Looking at the absolute numbers, the throughput is still lower than I
expected. I initially assumed that the distance computations would
dominate the execution time, since scanning a list means comparing
every stored vector against the query. Instead, most of the time is
spent \textbf{before} the scan even starts.

The reason is that reading a list from CTE currently requires
\textbf{copying the whole blob into a temporary buffer}. Only after
this copy has completed can the scanner begin computing distances. As
a result, every inverted list involved in a search is copied once
before it is processed, and for searches that touch many lists this
represents a large amount of memory traffic.

I confirmed that this copy, and not the distance computation, is the
bottleneck through three observations:

\begin{itemize}
\item The internal timers of the ChiMod show that the time spent
waiting for reads to complete is larger than the time spent actually
scanning.
\item Switching the distance computation from the scalar
implementation to the wider AVX-512 vector instructions changed the
throughput by almost nothing. If the search were limited by
arithmetic, faster arithmetic would help; the fact that it does not
means the search is limited by data movement.
\item In the regime where the index fits in memory, the warm
memory-mapped baseline is actually faster than the CTE version (about
70 QPS versus about 28 QPS at 24 gigabytes). The difference is exactly
the extra copy: with a warm memory-mapped file, the page that the
scanner reads is already the final location of the data, so no copy is
needed, whereas the CTE version copies the list before scanning it.
\end{itemize}

\subsection{Current Search Pipeline}

The current execution flow can be summarized as follows, where steps 2
to 5 overlap across different lists:

\begin{enumerate}
\item Determine the inverted lists that must be visited.
\item Issue an \texttt{AsyncGetBlob()} request for each list.
\item Wait until the blob has been copied into a temporary buffer.
\item Scan the copied data.
\item Discard the temporary buffer.
\end{enumerate}

Even when a list already resides in the RAM tier of CTE, the search
still performs an additional memory copy before any distance
computation can begin. Conceptually, the execution looks as follows.

\begin{verbatim}

CTE RAM

+------------------------+
| Inverted List          |
+------------------------+
           |
           | AsyncGetBlob()
           v

+------------------------+
| Temporary Buffer       |
+------------------------+
           |
           v

Distance Computation

\end{verbatim}

\subsection{Copies, Disk Reads, and Where the Time Goes}

It is worth being precise about which part of this cost is avoidable,
because the answer depends on where each list is stored.

When a list is in the \textbf{RAM tier}, the data is already in memory
and the copy into the temporary buffer is pure overhead: the same
bytes are duplicated for no reason other than the shape of the read
interface.

When a list is in the \textbf{NVMe tier}, the read is not pure
overhead: those bytes genuinely have to be brought from the solid
state disk into memory. This part is unavoidable with any storage
system.

This distinction matters for large indexes. When the index is several
times larger than RAM, most lists live on the NVMe tier, so a large
part of every search is genuine disk reading rather than wasteful
copying. Measuring the amount of data moved per search against the
elapsed time gives an effective throughput of roughly 0.7 to 0.8
gigabytes per second at 49 and 87 gigabytes, which is close to the
read bandwidth of the local NVMe device. In other words, at these
sizes the search is essentially reading the working set off the disk
on every query batch, and the copy is an additional cost on top of the
part that is already in RAM.

The general conclusion is that the remaining bottleneck is no longer
located in FAISS, nor in the operating system, but in the interface
used to access CTE data. The next section describes an extension to
CLIO that could remove the avoidable part of this cost.

%%%%%%%%%%%%%%%%%%%%%%%%%%%%%%%%%%%%%%%%%%%%%%%%%%%%%%%%%%%%%%%
\section{Possible Extension to CLIO}

The current implementation shows that CTE is a much better storage
abstraction for IVF search than relying on the operating system page
cache, because inverted lists are managed explicitly instead of
implicitly by the kernel. The experiments also show that the remaining
bottleneck is no longer storage placement, but the interface used to
retrieve data from CTE.

Currently, every call to \texttt{AsyncGetBlob()} copies the requested
blob into a temporary buffer before the application can access it. For
an application such as FAISS, where the data is only read and never
modified, this copy is unnecessary when the blob already resides in
the local RAM tier.

\subsection{Why This Cannot Be Done Today}

Before proposing a change, I looked at the CTE implementation to
understand why a module cannot simply read the data in place already.
There are three specific reasons:

\begin{itemize}
\item The bytes stored in the RAM tier live in memory that is private
to the runtime process, not in the shared-memory region that clients
normally address, so a client-side pointer cannot reach them directly.
\item There is a query that would report where a blob is physically
stored, but the fields that carry that information are currently
disabled, so no caller can resolve a blob to a concrete location.
\item Nothing prevents a blob from being moved between tiers or
reclaimed while it is being read. Reading its memory directly would
therefore risk accessing data that has already been freed or relocated.
\end{itemize}

The encouraging part is that, inside the runtime, a block of RAM keeps
a stable address for as long as it exists. The storage already behaves
in the way a direct read would need; what is missing is a safe way to
expose it.

\subsection{A Possible Improvement}

One possible extension would be to allow a ChiMod running inside the
same CLIO runtime to obtain a temporary read-only view of a blob that
is already stored in the local RAM tier. Conceptually, the workflow
would become:

\begin{verbatim}

Current

CTE RAM
    |
    | AsyncGetBlob()
    v
Temporary Buffer
    |
    v
Distance Computation


Proposed

CTE RAM
    |
    | Read-only view
    v
Distance Computation

\end{verbatim}

Instead of copying the contents into another buffer, the scanner would
process the bytes directly from their original location in the RAM
tier.

\subsection{Requirements}

From my understanding of the current implementation, this would
require three capabilities.

\begin{itemize}

\item \textbf{Resolve.} Given a blob identifier, determine where the
blob is currently stored.

\item \textbf{View.} If the blob is located in the local RAM tier,
obtain a temporary pointer to its contents.

\item \textbf{Pin.} Guarantee that the blob cannot be relocated or
reclaimed while it is being scanned.

\end{itemize}

These three operations compose naturally: resolve finds the data, view
exposes it, and pin keeps it safe for the duration of the scan. If a
blob is not resident in the RAM tier, the existing
\texttt{AsyncGetBlob()} mechanism could still be used as a fallback.

\subsection{Scope of the Improvement}

It is important to be honest about what this change would and would not
solve. It removes the copy for the part of the index that is resident
in the RAM tier, which is the avoidable overhead identified in the
experiments. It does not remove the cost of reading lists that live on
the NVMe tier, because those bytes still have to be brought from disk.

Therefore, this extension would bring the in-memory working set close
to the speed of a warm memory-mapped index, while keeping the
advantages that the page cache cannot offer: explicit tiers, no
page-cache thrashing, predictable behaviour, and placement decisions
made at the granularity of an inverted list. For indexes that are much
larger than RAM, the local disk bandwidth remains the ultimate limit,
and reducing it further would require a separate mechanism to read the
NVMe tier without a full copy.

%%%%%%%%%%%%%%%%%%%%%%%%%%%%%%%%%%%%%%%%%%%%%%%%%%%%%%%%%%%%%%%
\section{Conclusions}

The goal of this project was to investigate whether FAISS IVF search
could benefit from CLIO's explicit storage management.

The implementation successfully replaces \texttt{OnDiskInvertedLists}
with a CTE-based representation where each inverted list is stored as
an independent blob. This avoids the page-cache thrashing observed
with memory-mapped indexes once the dataset becomes much larger than
the available RAM, and it does so while producing exactly the same
search results as stock FAISS.

The evaluation also showed that, once the page cache is out of the
picture, the remaining bottleneck is the cost of copying every blob
into a temporary buffer before it can be processed. Part of this cost,
the part corresponding to lists already resident in the RAM tier, is
avoidable.

For this reason, I believe that providing a mechanism for obtaining
safe read-only views of blobs already resident in the local RAM tier,
built from the resolve, view and pin operations described above, would
significantly reduce the amount of memory movement performed during a
search. More generally, such functionality could also benefit other
near-data applications running inside CLIO, not only FAISS.

%%%%%%%%%%%%%%%%%%%%%%%%%%%%%%%%%%%%%%%%%%%%%%%%%%%%%%%%%%%%%%%
\section*{Future Work}

There are several directions that could be explored in the future.

\begin{itemize}

\item Implement the read-only view mechanism and measure the
performance improvement obtained by eliminating the avoidable copy.

\item Measure throughput under different RAM-tier capacities, to see
how the crossover between the RAM tier and the NVMe tier affects
performance.

\item Investigate a mechanism to read the NVMe tier without a full
copy, so that very large indexes are limited only by disk bandwidth.

\item Investigate whether the same read-only view technique could be
applied to other data-intensive applications running inside the CLIO
runtime.

\end{itemize}

\end{document}
